\section{Forecasting Architecture and Experimental Results}
\label{sec:forecast-results}

This section evaluates the proposed TSFM-Graph Adapter architecture purely as a
forecasting model, without incorporating any portfolio construction or trading
rules. The objective is to isolate the predictive value of the temporal branch,
the graph branch, and their fusion under a common rolling one-step-ahead
forecasting protocol.

\subsection{Model Variants}

We compare three model variants built on the same forecasting pipeline for the
target commodity \texttt{CO1 Comdty} (Brent crude oil):

\begin{enumerate}
    \item \textbf{Temporal-Only}: a baseline that uses only the frozen TimesFM
    backbone to extract temporal embeddings from the target series.
    \item \textbf{Graph-Only}: a baseline that removes the TimesFM temporal
    branch and predicts using only cross-asset graph information derived from
    handcrafted node features, a correlation-based adjacency matrix, and a graph
    attention network.
    \item \textbf{Total Fusion}: the proposed two-branch model that combines
    frozen TimesFM temporal embeddings with graph context through a gated fusion
    adapter.
\end{enumerate}

The graph branch operates over a multi-asset commodity universe and uses a
weighted correlation graph with a 60-day rolling window and absolute
correlations. Node representations are processed by a multi-layer graph
attention network, and the resulting target-node graph embedding is either used
directly (\emph{Graph-Only}) or injected into the temporal embedding sequence
through a gated late-fusion module (\emph{Total Fusion}).

\subsection{Experimental Setup}

All models were evaluated under the same out-of-sample rolling forecasting
protocol. At each test step, the model observed only information available up to
day $T$ and produced a one-step-ahead forecast for day $T+1$. The downstream
supervision target was defined in \texttt{log\_return} mode, while forecast
quality was assessed in terms of reconstructed price predictions. The test
window contained 810 one-step-ahead forecasts spanning the period from
\texttt{2023-03-28} to \texttt{2025-06-14}.

\subsection{Forecasting Results}

Table~\ref{tab:forecast-metrics} reports the out-of-sample forecasting
performance of the three architectures. We evaluate each model with MAPE,
$R^2$, MAE, MSE, RMSE, and directional accuracy.

\begin{table}[ht]
    \centering
    \caption{Out-of-sample forecasting performance without portfolio evaluation.}
    \label{tab:forecast-metrics}
    \begin{tabular}{lcccccc}
        \toprule
        \textbf{Model} & \textbf{MAPE (\%)} & \textbf{$R^2$} & \textbf{MAE} & \textbf{MSE} & \textbf{RMSE} & \textbf{DirAcc (\%)} \\
        \midrule
        Temporal-Only & 1.128 & 0.9688 & 0.8750 & 1.5382 & 1.2402 & 46.17 \\
        Graph-Only (cw60, abs) & 1.037 & 0.9702 & 0.8087 & 1.4721 & 1.2133 & 48.64 \\
        Total Fusion (cw60, abs) & 1.083 & 0.9697 & 0.8407 & 1.4980 & 1.2239 & 52.22 \\
        \bottomrule
    \end{tabular}
\end{table}

The results show a clear trade-off between absolute forecast error and
directional quality. The \textbf{Graph-Only} model achieves the lowest MAPE,
MAE, MSE, and RMSE, indicating that cross-asset graph structure alone carries
useful information for reconstructing the next-day price level. However, its
directional accuracy remains below 50\%, suggesting that good point forecasts do
not necessarily translate into reliable sign prediction.

The \textbf{Temporal-Only} baseline performs competitively in terms of absolute
error but yields the weakest directional accuracy at 46.17\%. This indicates
that the frozen temporal backbone alone is insufficient to consistently capture
the sign of the next-day move in the current setup.

Most importantly, the proposed \textbf{Total Fusion} model achieves the highest
directional accuracy at 52.22\%, outperforming both ablations in sign
prediction. Although its absolute error metrics are not the best among the three
models, the fusion architecture improves directional behavior relative to both
the temporal-only and graph-only baselines. This suggests that graph context is
most useful when it complements temporal embeddings rather than replacing them.

\subsection{Discussion}

These findings support the main architectural hypothesis of this study:
cross-sectional graph information and temporal sequence information play
different roles in one-step-ahead commodity forecasting. The temporal branch is
responsible for capturing asset-specific historical dynamics, while the graph
branch provides cross-asset context derived from correlation structure and node
features. When the two branches are used together, the model achieves the best
directional performance, even if the lowest pointwise forecast error is obtained
by the graph-only baseline.

From a forecasting perspective, this distinction is important. In commodity
applications, small improvements in directional accuracy may be more valuable
than marginal improvements in level error, especially when the downstream use
case depends on predicting the sign of short-horizon moves. Therefore, the
fusion model appears to offer the most balanced forecasting behavior among the
tested architectures.

\subsection{Qualitative Visualization}

Figure~\ref{fig:forecast-comparison} presents the out-of-sample forecast
comparison across all three architectures.

\begin{figure}[ht]
    \centering
    \includegraphics[width=0.9\linewidth]{forecast_comparison_graph_adapter.png}
    \caption{Rolling one-step-ahead forecast comparison for Temporal-Only,
    Graph-Only, and Total Fusion models on \texttt{CO1 Comdty}.}
    \label{fig:forecast-comparison}
\end{figure}

\subsection{Summary}

In summary, the empirical results indicate that:

\begin{itemize}
    \item the graph branch alone yields the lowest pointwise forecast error,
    \item the temporal branch alone is not sufficient for strong directional
    prediction,
    \item the fused architecture provides the best directional accuracy and the
    most balanced overall forecasting behavior.
\end{itemize}

For this reason, the \textbf{Total Fusion} model is the most promising variant
from a forecasting standpoint, while the ablations help clarify the distinct
contributions of temporal and cross-sectional graph information.
